\thaichapter{บทที่ 1}{บทนำ}\addcontentsline{toc}{chapter}{บทที่ 1 บทนำ}

\section*{1.1 ที่มาและความสำคัญ}\addcontentsline{toc}{section}{1.1 ที่มาและความสำคัญ}

เว็บไซต์ที่เปิดสู่อินเทอร์เน็ตถูกสแกนและทดลองเจาะระบบโดยโปรแกรมอัตโนมัติตลอดเวลา โดยเฉพาะเว็บไซต์ที่ใช้ WordPress ซึ่งเป็นระบบจัดการเนื้อหาที่ได้รับความนิยมมาก ผู้ดูแลเว็บไซต์จึงต้องเข้าใจทั้งการติดตั้งระบบและการป้องกันทุกชั้น ตั้งแต่เครือข่าย ระบบปฏิบัติการ เว็บเซิร์ฟเวอร์ ไปจนถึงตัวแอปพลิเคชัน

รายวิชานี้กำหนดให้นิสิตเช่าเซิร์ฟเวอร์บนคลาวด์และติดตั้ง WordPress ด้วยตนเอง แทนการใช้บริการสำเร็จรูป เพื่อให้ได้ฝึกตั้งค่าความปลอดภัยด้วยตนเองทุกขั้นตอน อาจารย์ผู้สอนแจ้งว่าจะใช้ AI ทดลองเจาะเว็บไซต์ของแต่ละกลุ่ม คณะผู้จัดทำจึงออกแบบการป้องกันหลายชั้น และพัฒนาระบบตรวจจับการโจมตีเพิ่มเติม เพื่อให้เห็นได้ว่าใครพยายามเข้าถึงระบบบ้าง และการป้องกันทำงานได้จริงหรือไม่

เนื้อหาของเว็บไซต์เน้นการประชาสัมพันธ์สาขาวิชาวิทยาการคอมพิวเตอร์ คณะวิทยาการสารสนเทศ มหาวิทยาลัยบูรพา โดยมีกลุ่มเป้าหมายหลักคือนักเรียนมัธยมศึกษาตอนปลายที่กำลังเลือกสาขาเรียน

\section*{1.2 วัตถุประสงค์}\addcontentsline{toc}{section}{1.2 วัตถุประสงค์}

\begin{enumerate}
\def\labelenumi{\arabic{enumi}.}
\tightlist
\item
  เพื่อเช่าเซิร์ฟเวอร์บนคลาวด์และติดตั้ง nginx, PHP, MySQL และ WordPress ด้วยตนเอง
\item
  เพื่อให้เว็บไซต์เข้าถึงได้ผ่านชื่อโดเมนด้วย HTTPS
\item
  เพื่อตั้งค่า SSH ให้ปิดการเข้าใช้ด้วย root ปิดการยืนยันตัวตนด้วยรหัสผ่าน และใช้กุญแจสาธารณะ (public key) เท่านั้น
\item
  เพื่อทำ firewall ที่รับการเชื่อมต่อเฉพาะจากหมายเลข IP ที่จดทะเบียนในประเทศไทย
\item
  เพื่อพัฒนาเว็บไซต์ประชาสัมพันธ์สาขาวิชาวิทยาการคอมพิวเตอร์ที่เหมาะกับนักเรียนมัธยมศึกษาตอนปลาย
\item
  เพื่อพัฒนาระบบตรวจจับการโจมตีจาก log พร้อม dashboard และการสรุปด้วย AI
\end{enumerate}

\section*{1.3 ขอบเขตของโครงงาน}\addcontentsline{toc}{section}{1.3 ขอบเขตของโครงงาน}

\begin{itemize}
\tightlist
\item
  เซิร์ฟเวอร์เสมือน (VM) 1 เครื่องบน Google Cloud ระบบปฏิบัติการ Ubuntu Server 24.04 LTS
\item
  โดเมน elephotelephus.site จดทะเบียนและจัดการ DNS ที่ Cloudflare
\item
  เว็บไซต์พัฒนาด้วย WordPress และธีมที่คณะผู้จัดทำเขียนขึ้นเอง
\item
  การป้องกันครอบคลุมเครือข่าย (firewall) การเข้าเครื่อง (SSH) เว็บเซิร์ฟเวอร์ WordPress และฐานข้อมูล
\item
  ระบบตรวจจับการโจมตีอ่าน log 4 แหล่ง ได้แก่ nginx, SSH, firewall และ fail2ban ย้อนหลัง 7 วัน
\item
  ไม่รวมการทำระบบสำรองหลายเครื่อง (high availability) และ CDN
\end{itemize}

\section*{1.4 ประโยชน์ที่คาดว่าจะได้รับ}\addcontentsline{toc}{section}{1.4 ประโยชน์ที่คาดว่าจะได้รับ}

\begin{enumerate}
\def\labelenumi{\arabic{enumi}.}
\tightlist
\item
  ได้เว็บไซต์ประชาสัมพันธ์สาขาที่ใช้งานได้จริง และทีมแก้ไขเนื้อหาได้เองโดยไม่ต้องแก้โค้ด
\item
  ได้ฝึกติดตั้งและดูแลเซิร์ฟเวอร์ Linux ตั้งแต่ต้นจนใช้งานจริง
\item
  เข้าใจหลักการป้องกันหลายชั้น (defense in depth) จากการลงมือทำ
\item
  มีเครื่องมือดูการโจมตีที่เกิดขึ้นจริง และใช้ประกอบการตัดสินใจปรับปรุงการป้องกัน
\end{enumerate}
